% Mathematical response source extracted from unchanged Cycle8 main proof.
\subsection{Inherited susceptibility and the exact target}

Put \(b=\beta^2\), \(a=1-b\), and
\[
 \sigma=\operatorname{diag}(1/2,1/4,1/4)^{\otimes k},
 \qquad-\log_2\sigma=kI+Q.
\]
The bounded physical displacement
\[
 G_\beta=Q-W_k^*QW_k,\qquad \norm{G_\beta}\le k
\]
equals \(W_k^*LW_k-L\) on the finite-coordinate core, so its definition
does not subtract possibly infinite depth expectations.
For a signed charge-preserving input \(X\), write
\[
 D_b(X)=\Tr_A[(X\otimes I)G_\beta].
\]
It commutes with \(L\). On initial depth at least \(k\), its input-charge
pinching has the inherited exact first-\(k\)-color representation
\[
 \mathcal G_k(b)=G_0+bG_1+R_b,\qquad
 G_0=(2Q-kI)\otimes I_F,\qquad F=(\mathbb C^2)^{\otimes k}.
\]
At most \(k\) initial colors are consumed; final untouched tails agree in
every contributing equal-charge pair. The representation is tensor identity
on the seed and remaining tail. It is a bulk identity, not a cutoff or a
replacement of the root.

For a prefix density \(\zeta\), define
\[
 H_\zeta=\Tr_F[(I\otimes\zeta)G_1]+2Q-kI,\qquad
 T_k(X)=\Tr_A[(X\otimes I)G_1].
\]
The unchanged Cycle7 theorem gives
\begin{align}
 v_k&=\min_{\zeta}\Tr(\sigma H_\zeta^2)
 =\frac{k(k-1)}8-\frac18
   \sum_{n=0}^{k-2}(k-1-n)\Prob(S_n=0),\label{eq:vk}\\
 H&=H_{\zeta_{\rm sym}},\qquad
 \zeta_{\rm sym}=\frac{P_{{\rm sym},k}}{k+1},
 \qquad T_k(\sigma H)=v_kI_F.\label{eq:scalar}
\end{align}
The walk \(S_n\) is simple symmetric and starts at zero. The minimum admits
all coherent and entangled prefix densities. The analytic density
\(\zeta_{\rm sym}\) is never physically prepared. Useful inherited bounds are
\begin{equation}\label{eq:inherited}
 \frac{k^2}{32}\le v_k\le\frac{k(k-1)}8,\quad
 v_k\ge\frac N8,\quad N=\binom{k}{2},\quad
 \norm H\le k(k-1),\quad
 \Tr\sigma H=\Tr\sigma QH=0.
\end{equation}
Also \(H=H^*\), \([H,Q]=[H,\sigma]=0\), and \(H\) is input-offdiagonal.

We need its actual block structure, not just its norm. Write
\(H=\sum_{i<j}H_{ij}\), where \(H_{ij}\) connects complete binary charge
patterns differing by \(01\leftrightarrow10\) at \(i,j\).
On an oriented block, the Cycle7 normal form is
\begin{equation}\label{eq:Hblocks}
 K=
 \begin{cases}
 P(I-\SWAP),&h>d,\\
 P/2,&h=d,\\
 0,&h<d.
 \end{cases}
\end{equation}
Here \(h\) is the unmatched earlier push count, \(d\) the unmatched
middle pop count, and \(P\) a color-wire permutation. The reversed block
is \(K^*\). The two swapped wires in the first case are distinct input
color qubits. These are inherited bulk coefficient identities.

The preceding sufficient window \(bk^4\le2^{-12}\) came from bounding
\(\norm{R_b}\le4k^5b^2\) before contraction. We instead control the exact
quantity \(\Tr_A[(\sigma H\otimes I)\mathcal G_k(b)]\).

\subsection{Exact response from paired coherent histories}

\begin{theorem}[All-order endpoint response]\label{thm:response}
For every \(k\ge2\) and \(0\le b\le1\), there is a Hermitian operator
polynomial \(F_k(b)\) on the entire prefix \(F\) such that
\begin{align}
 \Tr_A[(\sigma H\otimes I)\mathcal G_k(b)]
   &=b(1-b)F_k(b),\qquad F_k(0)=v_kI,\label{eq:Fidentity}\\
 \norm{F_k(b)-v_kI}
   &\le N[1-(1-b)^{k-2}]
    \le N(k-2)b.\label{eq:Fbound}
\end{align}
In particular
\begin{equation}\label{eq:contracted}
 \norm{\Tr_A[(\sigma H\otimes I)R_b]}
 \le b^2v_k+b(1-b)N[1-(1-b)^{k-2}]
 \le b^2[v_k+N(k-2)].
\end{equation}
Every inequality is an operator norm bound, not an expectation at one
selected prefix.
\end{theorem}

\subsubsection{Charge selection and the cancellation of output charge}

For a stay mask \(M\subseteq\{0,\ldots,k-1\}\), let \(U_M\) be the ordered
word with \(I\) at \(M\) and \(U\) elsewhere. Its coherent amplitude is
\[
 \alpha^{k-|M|}(i\beta)^{|M|}U_M.
\]
On a bulk computational charge pattern \(x\), the output charge at a site
is \(x_\ell\) for a stay and \(1-x_\ell\) for a move. Suppose two initial
patterns differ only at endpoints \(i,j\), with opposite charges there.
A common output charge word forces the two histories to make identical
choices at every spectator and complementary choices at both endpoints.

Fix the common spectator stay set \(S\). Its paired scalar weight is
\[
 w_S(b)=b^{|S|}a^{k-2-|S|},\qquad
 \sum_{S\subseteq[k]\setminus\{i,j\}}w_S(b)=1.
\]
The endpoints supply \(ba\). Use a ket stay mask
\(s=(s_i,s_j)\), with ket input charges \(01\), and let \(R_s\) denote
the compressed raw cross-history operator. The bra mask is complementary:
\[
\begin{array}{c|c|c|c}
s&\text{bra stay mask}&\text{phase sign}&
       \text{output endpoint charge}\\ \hline
00&11&-1&1\\
01&10&+1&2\\
10&01&+1&0\\
11&00&-1&1
\end{array}
\]
Each \(R_s\) has norm at most one. This retains all input color and final
bath overlaps.

Keep the spectator calls fixed to \(U\) or \(I\), with only the endpoints
equal to \(V_\beta\). This auxiliary word is unitary. Its offdiagonal
input block in \(W^*W\) vanishes, and the four terms have a common \(ba\)
factor. Hence, for \(0<b<1\), and then as a polynomial identity,
\begin{equation}\label{eq:unitarity}
 -R_{00}+R_{01}+R_{10}-R_{11}=0.
\end{equation}
The spectator output charge is constant on the selected charge-pattern
block. Equation~\eqref{eq:unitarity} cancels that entire contribution to
\(W^*QW\), and also cancels the constant one in the endpoint charge.
The surviving block in \(Q-W^*QW\) is
\[
 R_{10}-R_{01},
\]
a difference of two contractions. There is no factor \(k\) from
\(\norm Q\).

Let \(B_{ij}^S\) be this block, its adjoint and their direct sum over
the complete charge patterns. The pattern pairs are disjoint, so
\begin{equation}\label{eq:Bnorm}
 \norm{B_{ij}^S}\le2.
\end{equation}
Equivalently it is the \(i,j\) endpoint component of the first coefficient
of the ordered word with spectator calls \(S\) skipped, embedded on the
original input space. This is an analytic representation only; the actual
communication process still uses and counts all \(k\) calls.

Define
\begin{equation}\label{eq:A}
 A_{ij,S}=\Tr_A[(\sigma H_{ij}\otimes I)B_{ij}^S].
\end{equation}
The original \(k\)-call \(H_{ij}\) remains in this formula. It is not replaced
by an optimizer for a shorter word. Exact history pairing now gives
\begin{equation}\label{eq:Bernstein}
 F_k(b)=\sum_{i<j}\ \sum_{S\subseteq[k]\setminus\{i,j\}}
          b^{|S|}a^{k-2-|S|}A_{ij,S}.
\end{equation}
Each \(A_{ij,S}\) is Hermitian by input partial-trace cyclicity. No
commutation of \(H_{ij}\) with \(B_{ij}^S\) is required.

The positive spectator weights arise after matching the two coherent
histories. Equation~\eqref{eq:Bernstein} does not make the physical
device a classical independent-error channel.

\subsubsection{The weighted trace norm that makes the estimate uniform}

\begin{lemma}[Endpoint trace budget]\label{lem:trace}
For the actual inherited \(H\),
\[
 \norm{\sigma H_{ij}}_1\le\frac14,\qquad
 \norm{A_{ij,S}}\le\frac12.
\]
\end{lemma}
\begin{proof}
Conditional on a complete charge pattern with \(q\) colors, the input
color weight is \(I/2^q\). In the first case of \eqref{eq:Hblocks},
\[
 |P(I-\SWAP)|=I-\SWAP,\qquad
 2^{-q}\Tr(I-\SWAP)=\frac12.
\]
The swap has normalized trace \(1/2\) on two qubits. In the critical
case \(|P/2|=I/2\), with the same normalized trace. The final case is
zero. The two orientations have identical singular-value sums, while
their combined endpoint-charge probability under \(\sigma\) is \(1/2\).
Since \(H_{ij}\) commutes with \(\sigma\),
\[
 \norm{\sigma H_{ij}}_1
 =\Tr(\sigma|H_{ij}|)\le\frac12\frac12=\frac14.
\]

For a Hermitian input operator \(X\) and bounded joint operator \(B\),
\[
 \norm{\Tr_A[(X\otimes I)B]}\le\norm X_1\norm B.
\]
Diagonalize \(X\), bound the norm of each corresponding input diagonal
block of \(B\), and sum with the absolute eigenvalues. This remains valid
with an infinite bath. Apply it to \eqref{eq:A} and \eqref{eq:Bnorm}.
\end{proof}

The empty-mask certificate is
\(\sum_{i<j}A_{ij,\varnothing}=v_kI\), by \eqref{eq:scalar}.
For each pair, the nonempty-mask mass is \(1-a^{k-2}\).
Lemma~\ref{lem:trace} and the triangle inequality give
\[
 \left\|\sum_Sw_S A_{ij,S}-A_{ij,\varnothing}\right\|
 \le1-a^{k-2}.
\]
Sum over \(N\) pairs and use Bernoulli's inequality to prove
\eqref{eq:Fbound}. Finally
\[
 \Tr_A[(\sigma H\otimes I)R_b]
 =ba(F_k(b)-v_kI)-b^2v_kI,
\]
because \(H\) is input-offdiagonal. This proves
Theorem~\ref{thm:response}.

\subsubsection{Reusable interface and its actual premises}

For any Hermitian pair-supported \(J=\sum J_{ij}\) commuting with
\(\sigma\), the same proof gives \eqref{eq:Bernstein} with \(J\) in place
of \(H\), and
\[
 \norm{F_J(b)-F_J(0)}
 \le4\!\left(\sum_{i<j}\norm{\sigma J_{ij}}_1\right)
      [1-a^{k-2}].
\]
A known nonscalar first operator certificate is allowed. The scalar
certificate and quarter trace budgets above belong to the actual Cycle7
optimizer. General total-charge preservation does not imply pair support,
and the root does not satisfy the bulk charge-flip premise. Concrete
counterexamples are preserved in the unchanged Cycle8 source \cite{Cycle8}.

